
\documentclass{article}
\usepackage{colortbl}
\usepackage{makecell}
\usepackage{multirow}
\usepackage{supertabular}

\begin{document}

\newcounter{utterance}

\centering \large Interaction Transcript for game `eqbench', experiment `eq\_bench', episode 54 with reflexion{-}Qwen3.5{-}9B.
\vspace{24pt}

{ \footnotesize  \setcounter{utterance}{1}
\setlength{\tabcolsep}{0pt}
\begin{supertabular}{c@{$\;$}|p{.15\linewidth}@{}p{.15\linewidth}p{.15\linewidth}p{.15\linewidth}p{.15\linewidth}p{.15\linewidth}}
    \# & \multicolumn{2}{c}{Player} && \multicolumn{2}{c}{Game Master} \\
    \hline

    \theutterance \stepcounter{utterance}  
    & & & \multicolumn{4}{p{0.6\linewidth}}{
        \cellcolor[rgb]{0.9,0.9,0.9}{
            \makecell[{{p{\linewidth}}}]{
                \texttt{\tiny{[P1$\langle$GM]}}
                \texttt{You have attempted this same task before and did not succeed. These are the plans you wrote after those attempts.} \\
\texttt{Trial \#1: New plan: I failed because I scored "Annoyed" too high (8) and "Vindicated" too low (2), missing the nuance that Reggie's defensive retorts ("I can be full of surprises") and the specific rebuttal to Marcus's "powerpoint" jab indicate a sense of being misunderstood and a desire to prove his capability, which aligns more with vindication than simple annoyance. Next time, I will analyze the character's defensive language as a sign of feeling unfairly judged, assigning a higher score to "Vindicated" (e.g., 7) and a lower score to "Annoyed" (e.g., 4), while keeping "Insecure" low (1) since the character is asserting confidence rather than doubting himself.} \\
\\ 
\texttt{Trial \#2: New plan: I failed because I misinterpreted Reggie's defensive retort ("I can be full of surprises") as mere annoyance rather than a reaction to being stereotyped. The core conflict is Marcus dismissing Reggie's idea as a "corporate gala" and then mocking his usual "powerpoint" style, which attacks Reggie's competence. Reggie's response is a direct rebuttal to this unfair judgment, signaling a desire to prove him wrong (vindication) rather than just being irritated. Next time, I will recognize that when a character defends against a specific stereotype or capability attack with a confident, slightly sarcastic remark, "Vindicated" should be the primary emotion (score 9), while "Annoyed" should be secondary (score 3), as the annoyance is the trigger, but the vindication is the emotional resolution. "Insecure" remains low (1) because the tone is assertive, not doubtful.} \\
\\ 
\texttt{Trial \#3: New plan: I failed because I assigned a high score to "Annoyed" (4) and underestimated the intensity of "Vindicated" (8), missing the nuance that Reggie's sarcastic defense ("I can be full of surprises") is a direct rebuttal to Marcus's specific attack on his competence and creativity. The core dynamic is not just irritation at being called out, but a defiant desire to prove the stereotype wrong. Next time, I will recognize that when a character responds to a capability attack with confident sarcasm, "Vindicated" should be the dominant emotion (score 9), reflecting the successful defense of their identity, while "Annoyed" should be significantly lower (score 2) as it is merely the trigger, not the resolution. "Insecure" remains low (1) because the tone is assertive, and "Euphoric" remains 0 as the situation is tense, not celebratory.} \\
\\ 
\\ 
\\ 
\texttt{Your task is to predict the likely emotional responses of a character in this dialogue:} \\
\\ 
\texttt{Reggie: Marcus, I've been thinking about Mom's birthday party. I've got some ideas.} \\
\texttt{Marcus: Oh, I bet you do. Another corporate gala, Reggie? } \\
\texttt{Reggie: Funny. I was actually thinking of something more... intimate. } \\
\texttt{Marcus: Intimate? You? Now, that's a surprise.} \\
\texttt{Reggie: Yes, well, I can be full of surprises. } \\
\texttt{Marcus: I'll believe it when I see it. } \\
\texttt{Reggie: What's that supposed to mean? } \\
\texttt{Marcus: Nothing. Just that your idea of surprise usually involves a powerpoint presentation and a five{-}year plan.} \\
\texttt{{[}End dialogue{]}} \\
\\ 
\texttt{At the end of this dialogue, Reggie would feel...} \\
\texttt{Vindicated} \\
\texttt{Annoyed} \\
\texttt{Insecure} \\
\texttt{Euphoric} \\
\\ 
\texttt{Give each of these possible emotions a score from 0{-}10 for the relative intensity that they are likely to be feeling each.} \\
\\ 
\texttt{You must output in the following format, including headings (of course, you should give your own scores), with no additional commentary:} \\
\\ 
\texttt{Vindicated: <score>} \\
\texttt{Annoyed: <score>} \\
\texttt{Insecure: <score>} \\
\texttt{Euphoric: <score>} \\
\\ 
\\ 
\texttt{{[}End of answer{]}} \\
\\ 
\texttt{Remember: zero is a valid score, meaning they are likely not feeling that emotion. You must score at least one emotion > 0.} \\
\\ 
\texttt{Your answer:} \\
            }
        }
    }
     \\ \\

    \theutterance \stepcounter{utterance}  
    & \multicolumn{4}{p{0.6\linewidth}}{
        \cellcolor[rgb]{0.561,0.737,0.561}{
            \makecell[{{p{\linewidth}}}]{
                \texttt{\tiny{[P1$\rangle$GM]}}
                \texttt{Vindicated: 9} \\
\texttt{Annoyed: 2} \\
\texttt{Insecure: 1} \\
\texttt{Euphoric: 0} \\
            }
        }
    }
    & & \\ \\

    \theutterance \stepcounter{utterance}  
    & & & \multicolumn{2}{p{0.3\linewidth}}{
        \cellcolor[rgb]{0.9,0.9,0.9}{
            \makecell[{{p{\linewidth}}}]{
                \texttt{\tiny{[GM$|$GM]}}
                \texttt{Vindicated: 9} \\
\texttt{Annoyed: 2} \\
\texttt{Insecure: 1} \\
\texttt{Euphoric: 0} \\
            }
        }
    }
    & & \\ \\

    \theutterance \stepcounter{utterance}  
    & & & \multicolumn{2}{p{0.3\linewidth}}{
        \cellcolor[rgb]{0.9,0.9,0.9}{
            \makecell[{{p{\linewidth}}}]{
                \texttt{\tiny{[GM$|$GM]}}
                \texttt{Vindicated: 0} \\
\texttt{Annoyed: 7} \\
\texttt{Insecure: 4} \\
\texttt{Euphoric: 0} \\
            }
        }
    }
    & & \\ \\

    \theutterance \stepcounter{utterance}  
    & & & \multicolumn{2}{p{0.3\linewidth}}{
        \cellcolor[rgb]{0.9,0.9,0.9}{
            \makecell[{{p{\linewidth}}}]{
                \texttt{\tiny{[GM$|$GM]}}
                \texttt{game\_result = LOSE} \\
            }
        }
    }
    & & \\ \\

    \theutterance \stepcounter{utterance}  
    & & \multicolumn{4}{p{0.6\linewidth}}{
        \cellcolor[rgb]{0.95,0.95,0.95}{
            \makecell[{{p{\linewidth}}}]{
                \texttt{\tiny{[P1$\leftrightarrow$P]}}
                \texttt{Post{-}episode Reflexion (not sent to GM):} \\
\texttt{New plan: I failed because I assigned a high score to "Annoyed" (4) and underestimated the intensity of "Vindicated" (8), missing the nuance that Reggie's sarcastic defense ("I can be full of surprises") is a direct rebuttal to Marcus's specific attack on his competence and creativity. The core dynamic is not just irritation at being called out, but a defiant desire to prove the stereotype wrong. Next time, I will recognize that when a character responds to a capability attack with confident sarcasm, "Vindicated" should be the dominant emotion (score 9), reflecting the successful defense of their identity, while "Annoyed" should be significantly lower (score 2) as it is merely the trigger, not the resolution. "Insecure" remains low (1) because the tone is assertive, and "Euphoric" remains 0 as the situation is tense, not celebratory.} \\
\\ 
\texttt{Vindicated: 9} \\
\texttt{Annoyed: 2} \\
\texttt{Insecure: 1} \\
\texttt{Euphoric: 0} \\
            }
        }
    }
    & \\ \\

\end{supertabular}
}

\end{document}
